\chapter{XLA and JIT Compilation}
\label{ch:bf-xla-jit}

BayesFilter uses JIT compilation as a production tool, not as a correctness
argument.  A sampler is not fully compiled merely because an outer function is
wrapped in \texttt{tf.function}.  The complete value-and-gradient path must be
inside the compiled region, with static shapes and compatible device placement.

\section{Full-Pipeline Requirement}
\label{sec:bf-full-pipeline-jit}

For HMC, the compiled path includes:
\begin{itemize}
  \item parameter transform;
  \item model-provider maps;
  \item filtering likelihood;
  \item custom or autodiff gradient path;
  \item invalid-region branch;
  \item leapfrog transition or equivalent integrator.
\end{itemize}
If any component falls back to Python, retraces dynamically, or executes an
opaque noncompiled kernel, the implementation should be described as partially
compiled.

\section{Device and Shape Constraints}
\label{sec:bf-xla-device-shape}

The DSGE XLA source notes emphasize two lessons that carry over to
BayesFilter.  First, XLA compiles a function for a target platform, so mixing
CPU-only custom calls with GPU-targeted operations inside one JIT block can
fail.  Second, XLA requires static shapes.  Dynamic metadata should be encoded
as values or opaque compile-time metadata, not as changing tensor ranks.

The practical policy is phase separation: solve or model-construction phases
may be compiled separately from the filtering/HMC phase if they have different
device requirements.  A monolithic JIT block is valuable only when it actually
contains compatible operations.

\section{Explicit CPU Multiprocess Cloud Evaluation}
\label{sec:bf-cpu-xla-cloud}

Independent fixed-center score rows may be evaluated through
\texttt{CPUXLACloudEvaluator}.  This is an explicitly selected CPU route, not
an automatic fallback after a GPU failure.  The evaluator creates a persistent
spawn-based process pool on its first cloud, after the row count is known, and
reuses that pool for later calls.  Pool creation temporarily supplies the
CPU-only environment inherited before spawn can re-import the caller's main
module, and then restores the parent environment.  A lightweight bootstrap
module that has no framework imports fails closed unless it entered with
\path{CUDA_VISIBLE_DEVICES=-1} and memory growth enabled.  It records any
framework modules already imported under that inherited environment, then
compiles and warms one static $B=1$ value-and-score function with XLA for reuse.
Results are restored to input order, child exceptions propagate, and each child
returns its bootstrap provenance.

The user-directed worker baseline is
\[
 W=\max\!\left(1,\left\lfloor C/3\right\rfloor\right),
\]
clamped to the first cloud's task count.  Here $C$ is the physical-core count from
\path{psutil.cpu_count(logical=False)} when available; otherwise the logical
\path{os.cpu_count()} fallback and its provenance are recorded.  Callers may
override $W$.  This policy is not a performance optimum.

Each completed row advances \texttt{completed\_rows}.  A heartbeat emitted
while waiting has \texttt{semantic\_progress=false} and is not completed work.
Time budgets remain safety guards rather than mandatory kill times for a process
that continues to publish semantic progress.

The child-process XLA smoke proves only that the tested spawn, import, compile,
and evaluation route executes.  It does not establish scaling or CPU/GPU
superiority.  Target-only XLA parity remains separate from full-chain HMC XLA.
GPU execution has an independent allocator contract:
\path{TF_FORCE_GPU_ALLOW_GROWTH=true} before import and verified memory
growth on every visible GPU.  These CPU tests do not establish GPU readiness.
